\documentclass[10pt,a4paper]{article}
\usepackage[top=1.65cm,bottom=1.65cm,left=1.85cm,right=1.85cm,footskip=0.8cm]{geometry}
\usepackage[T1]{fontenc}
\usepackage[utf8]{inputenc}
\usepackage[default,type1]{sourcesans}
\input{glyphtounicode}
\pdfgentounicode=1
\usepackage[dvipsnames]{xcolor}
\definecolor{primaryColor}{RGB}{25,64,89}
\usepackage{titlesec,enumitem,needspace,fancyhdr,lastpage}
\usepackage{microtype}
\usepackage[pdftitle={Ali Naderi Parizi - CV - AI and Systems Engineering},pdfauthor={Ali Naderi Parizi},colorlinks=true,urlcolor=primaryColor]{hyperref}
\urlstyle{same}
\setlength{\parindent}{0pt}
\setlength{\parskip}{2pt}
\setcounter{secnumdepth}{0}
\titleformat{\section}{\large\bfseries\color{primaryColor}}{}{0pt}{}[\vspace{3pt}\titlerule]
\titlespacing*{\section}{0pt}{9pt}{5pt}
\setlist[itemize]{leftmargin=1.25em,topsep=2pt,itemsep=1pt,parsep=0pt,partopsep=0pt}
\pagestyle{fancy}
\fancyhf{}
\renewcommand{\headrulewidth}{0pt}
\fancyfoot[L]{\footnotesize\color{gray}Ali Naderi Parizi\quad |\quad Curriculum Vitae}
\fancyfoot[R]{\footnotesize\color{gray}\thepage\ /\ \pageref*{LastPage}}
\newcommand{\entry}[3]{\needspace{4\baselineskip}\par\addvspace{4pt}\textbf{#1}\hfill {\small #2}\par\nopagebreak{\itshape #3}\par\nopagebreak}
\newcommand{\project}[2]{\needspace{4\baselineskip}\par\addvspace{4pt}\textbf{#1}\par\nopagebreak{\small #2}\par\nopagebreak}

\begin{document}
\begin{center}
{\fontsize{25}{29}\selectfont\bfseries\color{primaryColor}Ali Naderi Parizi}\par
\vspace{3pt}
{\large AI \& Systems Engineering \enspace $\cdot$ \enspace Technical Leadership}\par
\vspace{4pt}
\href{mailto:me@alinaderiparizi.com}{me@alinaderiparizi.com}\enspace |\enspace
\href{https://alinaderiparizi.com}{alinaderiparizi.com}\enspace |\enspace
\href{https://github.com/mralinp}{github.com/mralinp}\\[2pt]
\href{https://www.linkedin.com/in/ali-naderi-parizi-5a0a74107/}{LinkedIn}\enspace |\enspace
\href{https://scholar.google.com/citations?user=2zycrawAAAAJ}{Google Scholar}\enspace |\enspace Tehran, Iran
\end{center}
\vspace{-6pt}

\section{Profile}
Hands-on CTO at a small AI startup, systems engineer, and machine learning researcher. I take systems from exploratory R\&D and prototypes through production: AI applications and infrastructure, passwordless authentication, and embedded Linux. My research combines geometric shape analysis and deep learning for 3D breast ultrasound, with a journal publication and open-source research software. I connect work across models, backend services, operating systems, and hardware, and lead engineering through architecture, implementation, and mentoring.

\section{Professional Experience}
\entry{Chief Technology Officer, SciNext}{2025--present}{Small AI startup \enspace |\enspace Tehran, Iran}
\begin{itemize}
\item Own architecture and hands-on delivery across the production platform and AI applications; translate product needs into deployed systems and mentor engineers and early-stage AI teams.
\item Designed and built Kubernetes (k3s) infrastructure with Argo CD, CI pipelines, a private registry, and automated service onboarding and deployment.
\item Integrated Prometheus, Grafana, Loki, and Tempo; implemented organization-level isolation and per-application data credentials.
\item Developed AI products including LLM-driven data extraction and a Persian-first chemical compound database with a retrieval-augmented assistant.
\end{itemize}
\entry{Backend Software Engineer, Symmio}{2026}{Remote}
\begin{itemize}
\item Improved microservice reliability and performance; contributed to anomaly detection using service metrics and observability tooling.
\end{itemize}
\entry{Software Engineer (AI), CyberDesign Nordic AB}{2024--2025}{Stockholm, Sweden}
\begin{itemize}
\item Developed machine learning systems for indoor positioning from mobile sensor data, with backend APIs and data pipelines for training and deployment.
\end{itemize}
\entry{Software Engineer; Engineering Lead, IDmelon Technologies}{2022--2024}{Vancouver, Canada}
\begin{itemize}
\item Took FIDO2 passwordless-authentication components from initial R\&D and prototypes into production.
\item Built Windows local-account sign-in with security keys: a C++ credential provider, background authentication service, and local FIDO server.
\item Led component design and architecture, mentored junior engineers, and worked with QA to diagnose and resolve production defects.
\end{itemize}
\entry{R\&D Intern; Junior Embedded Engineer, Sico Systems Intelligizer}{2017--2019}{Shiraz, Iran}
Built FPGA and microcontroller systems, a DTMF-based GSM/GPRS control device, and an Allwinner F1C100S Linux board with analog video processing.

\section{Education}
\textbf{MSc, Artificial Intelligence and Robotics} \hfill 2020--2023\\
Iran University of Science and Technology, Tehran. Thesis: breast lesion classification in 3D automated breast ultrasound; advised by Dr. Mohsen Soryani and Dr. Ehsan Kozegar.\par
\textbf{BSc, Computer Engineering, Shiraz University} \hfill 2015--2020\\
Graduation project: a Linux-based system-on-board for analog video acquisition and embedded image processing.

\newpage
\setlength{\parskip}{1pt}
\setlist[itemize]{leftmargin=1.25em,topsep=1pt,itemsep=0pt,parsep=0pt,partopsep=0pt}
\titlespacing*{\section}{0pt}{6pt}{3pt}
\section{Research \& Publication}
\entry{3D Breast Ultrasound and Spectral Shape Analysis}{2020--present}{MSc research and continuing collaboration, Iran Image Processing Lab}
\begin{itemize}
\item Study lesion geometry for breast-mass classification; reproduced radiomics and hybrid CNN methods on TDSC-ABUS under a documented protocol.
\item Develop lesion meshing, generalized surface eigenmodes, and heat/wave kernel signatures; check spectra, scaling, and rigid motion. Classification and robustness studies are ongoing.
\item Maintain the \href{https://github.com/mralinp/abus-classification}{open research code and experimental protocol} and the \href{https://pypi.org/project/tdsc-abus2023-pytorch/}{TDSC-ABUS2023 PyTorch dataset toolkit}, for volumetric data, tumor crops, and memory-mapped caching.
\end{itemize}
S. Barekatrezaei, \textbf{A. Naderiparizi}, E. Kozegar, J. Ghofrani, and M. Soryani.
``Breast mass classification in 3D ABUS based on Laplace-Beltrami spectra and dual path CNN.''
\textit{Expert Systems with Applications}, vol. 299, article 129973, 2026.
\href{https://doi.org/10.1016/j.eswa.2025.129973}{doi:10.1016/j.eswa.2025.129973}.

\section{Selected Projects}
\project{Native OpenWrt for the ZLT X28 5G Router}{Embedded Linux, firmware reverse engineering \enspace |\enspace \href{https://github.com/mralinp/zlt-x28-openwrt}{GitHub}}
\begin{itemize}
\item Ported OpenWrt userspace to a locked-down MediaTek MT6890/T750 router; integrated its vendor kernel and cellular stack with netifd, LuCI, and standard networking.
\item Built guarded USB recovery, verified backups of 51 NAND partitions, and reproducible images with stock A/B fallback. Native firmware carries 5G traffic on one tested unit; broader validation is ongoing.
\end{itemize}
\project{Ez-Crawl: LLM-Driven Web Data Pipelines}{AI applications and developer tooling \enspace |\enspace \href{https://ezcrawl.ir}{ezcrawl.ir}}
\begin{itemize}
\item Built visual web-extraction pipelines using local or hosted LLMs, with branching, versioned workflows, and run history.
\item Implemented an MCP interface for agents to construct, validate, execute, and inspect extraction pipelines.
\end{itemize}
\project{MQute}{Async Python MQTT framework \enspace |\enspace \href{https://github.com/mralinp/MQute}{GitHub}}
FastAPI-style routing, typed payloads, dependency injection, middleware, and broker-free testing for MQTT services.
\project{JoME and Persian Speech Synthesis}{Embedded control and generative speech \enspace |\enspace \href{https://jome-farmer.ir}{JoME} \enspace |\enspace \href{https://tts.jome-farmer.ir}{TTS demo}}
Built an ESP32/FreeRTOS irrigation prototype with scheduling, safety interlocks, flow sensing, local control, and a FastAPI application; cloud integration is ongoing. Separately developed Persian TTS and voice-cloning pipelines, including data preparation and model adaptation.

\section{Teaching}
\begin{itemize}
\item \textbf{Computer Vision TA, IUST (2022):} supervised graduate labs and designed a 3D tumor-mesh classification project focused on lesion-surface detail.
\item \textbf{Data Structures and Algorithm Design TA, Shiraz University (2017):} designed assignments and projects, ran quizzes, and graded submissions.
\item \textbf{Elected ACM ICPC representative, Shiraz University (2016--2017):} organized contests, taught problem solving, and mentored competitors; competitive programming TA (2016).
\item \textbf{Informatics Olympiad Instructor, Dastgheyb High School, Sampad (2016):} taught algorithm design and competitive programming to high-school students.
\end{itemize}

\section{Academic Service}
\begin{itemize}
\item \textbf{Peer reviewer, \textit{Journal of Imaging Informatics in Medicine} (2026):} evaluated breast-imaging classification research and experimental methodology.
\item \textbf{Judge, AI Research Challenge, SID and IranDoc (2026):} evaluated scientific-literature analysis systems, including knowledge graphs, topic clustering, and research-trend analysis.
\end{itemize}

\section{Technical Skills}
\textbf{AI and research:} Python, PyTorch, Lightning, Transformers, NumPy, SciPy, 3D imaging, spectral geometry.\\
\textbf{Applications:} TypeScript, React, FastAPI, PostgreSQL, MongoDB, MQTT, LLM tooling.\\
\textbf{Systems:} C/C++, Linux, FIDO2, Windows credential providers, embedded Linux, ESP32, FreeRTOS, FPGA.\\
\textbf{Infrastructure:} Docker, Kubernetes, Argo CD, CI/CD, Prometheus, Grafana, Loki, Tempo.
\vfill
{\footnotesize\color{gray}Updated November 2026. Project links provide code, technical write-ups, and validation details.}
\end{document}
